% GENERATED FILE. Edit canonical lessons, then run npm run generate:books.
% canonical-source-hash: fnv1a64:846e4146020d7de6
% canonical-lessons: ZH-C260-nengli, ZH-C260-benshi, ZH-C260-gaoshou, ZH-C260-jishu, ZH-C260-jinbu

\chapter{What I Can Do}
\label{ch:zh-what-i-can-do}
\hlchaptermodality{260}

\begin{chapteropening}
\textbf{By the end of this chapter:} \emph{I can talk about ability, what someone can do, an expert, technique and making progress.}

Complete the last lesson of chapter 260: I can talk about ability, what someone can do, an expert, technique and making progress.
\end{chapteropening}

\section[nénglì]{nénglì — ability}

\label{lesson:ZH-C260-nengli}



\begin{quote}
Before the new one: say the Mandarin for a journalist, then the Mandarin for an electrician.
\end{quote}

\subsection*{You'll want to know: nénglì}
\emph{nénglì} — \textquotedblleft{}ability\textquotedblright{}.

\begin{quote}
\emph{tā hěn yǒu nénglì} — He's very capable.
\end{quote}

Its characters come later; for now it is a word for the ear.

\begin{cousinweb}
\emph{nénglì} is built from \zh{能} 'able' and \emph{lì} 'strength', both inherited from Middle Chinese, \emph{nong} and \emph{lik} (Guangyun, in Baxter's transcription), which glosses the first 'skilled' and the second 'sinew'.
\end{cousinweb}

\begin{grammarlens}[title={What you've built}]
One more word for what I have and can do.
\end{grammarlens}

\subsection*{Your turn}
\begin{itemize}
\raggedright
  \item \emph{Say it:} \emph{nénglì}
  \item \emph{Say it:} \emph{nénglì}, once more
  \item \emph{Say it:} \emph{tā hěn yǒu nénglì} — He's very capable.
  \item \emph{Recall:} say the Mandarin for a journalist, then the Mandarin for an electrician, then say \emph{nénglì} again
\end{itemize}

\begin{tcolorbox}[breakable,colback=teal!4,colframe=teal!35!black,title={Before you move on}]
What does nénglì mean? (\textquotedblleft{}Ability\textquotedblright{}.) Say it once more.
\end{tcolorbox}
\section[\texorpdfstring{\zh{本事} (běnshi)}{本事 (běnshi)}]{\zh{本事} (běnshi) — skill; what someone is capable of}

\label{lesson:ZH-C260-benshi}



\begin{quote}
Before the new one: say the Mandarin for an electrician, then the Mandarin for ability.
\end{quote}

\subsection*{You'll want to know: \zh{本事}}
\textbf{\zh{本事}} — \emph{běnshi} — \textquotedblleft{}skill; what someone is capable of\textquotedblright{}.

\begin{quote}
\textbf{\zh{这个人很有本事}} \emph{zhège rén hěn yǒu běnshi} — This man can turn his hand to anything.
\end{quote}

\begin{cousinweb}
\textbf{\zh{本事}} (\emph{běnshi}) is built from \zh{本} 'the root' and \zh{事} 'affairs', both inherited from Middle Chinese, \emph{pwonX} and \emph{dzriH} (Guangyun, in Baxter's transcription), which glosses the first 'the root' and 'the beginning'. In the word the second loses its tone.
\end{cousinweb}

\begin{grammarlens}[title={What you've built}]
One more word for what I have and can do.
\end{grammarlens}

\subsection*{Your turn}
\begin{itemize}
\raggedright
  \item \emph{Say it:} \emph{běnshi}
  \item \emph{Say it:} \emph{běnshi}, once more
  \item \emph{Say it:} \emph{zhège rén hěn yǒu běnshi} — This man can turn his hand to anything.
  \item \emph{Recall:} say the Mandarin for an electrician, then the Mandarin for ability, then say \emph{běnshi} again
\end{itemize}

\begin{tcolorbox}[breakable,colback=teal!4,colframe=teal!35!black,title={Before you move on}]
What does \zh{本事} mean? (\textquotedblleft{}Skill; what someone is capable of\textquotedblright{}.) Say it once more.
\end{tcolorbox}
\section[\texorpdfstring{\zh{高手} (gāoshǒu)}{高手 (gāoshǒu)}]{\zh{高手} (gāoshǒu) — expert; master (at something)}

\label{lesson:ZH-C260-gaoshou}



\begin{quote}
Before the new one: say the Mandarin for ability, then the Mandarin for what someone is capable of.
\end{quote}

\subsection*{You'll want to know: \zh{高手}}
\textbf{\zh{高手}} — \emph{gāoshǒu} — \textquotedblleft{}expert; master (at something)\textquotedblright{}.

\begin{quote}
\textbf{\zh{做饭他是高手}} \emph{zuòfàn tā shì gāoshǒu} — When it comes to cooking, he's an expert.
\end{quote}

\begin{cousinweb}
\textbf{\zh{高手}} (\emph{gāoshǒu}) is built from \zh{高} 'high' and \zh{手} 'a hand', both inherited from Middle Chinese, \emph{kaw} and \emph{syuwX} (Guangyun, in Baxter's transcription): a hand of a high level.
\end{cousinweb}

\begin{grammarlens}[title={What you've built}]
One more word for what I have and can do.
\end{grammarlens}

\subsection*{Your turn}
\begin{itemize}
\raggedright
  \item \emph{Say it:} \emph{gāoshǒu}
  \item \emph{Say it:} \emph{gāoshǒu}, once more
  \item \emph{Say it:} \emph{zuòfàn tā shì gāoshǒu} — When it comes to cooking, he's an expert.
  \item \emph{Recall:} say the Mandarin for ability, then the Mandarin for what someone is capable of, then say \emph{gāoshǒu} again
\end{itemize}

\begin{tcolorbox}[breakable,colback=teal!4,colframe=teal!35!black,title={Before you move on}]
What does \zh{高手} mean? (\textquotedblleft{}Expert; master (at something)\textquotedblright{}.) Say it once more.
\end{tcolorbox}
\section[jìshù]{jìshù — skill; technique}

\label{lesson:ZH-C260-jishu}



\begin{quote}
Before the new one: say the Mandarin for what someone is capable of, then the Mandarin for an expert.
\end{quote}

\subsection*{You'll want to know: jìshù}
\emph{jìshù} — \textquotedblleft{}skill; technique\textquotedblright{}.

\begin{quote}
\emph{tā de jìshù hěn hǎo} — He's very skilled.
\end{quote}

Its characters come later; for now it is a word for the ear.

\begin{cousinweb}
\emph{jìshù} is built from \emph{jì} 'a craft' and \emph{shù}, both inherited from Middle Chinese, \emph{gjeX} and \emph{zywit} (Guangyun, in Baxter's transcription), which glosses the first 'skill' and the second with this very word. Middle Chinese \emph{zywit} would regularly give a rising tone, shú; the falling tone of shù is irregular.
\end{cousinweb}

\begin{grammarlens}[title={What you've built}]
One more word for what I have and can do.
\end{grammarlens}

\subsection*{Your turn}
\begin{itemize}
\raggedright
  \item \emph{Say it:} \emph{jìshù}
  \item \emph{Say it:} \emph{jìshù}, once more
  \item \emph{Say it:} \emph{tā de jìshù hěn hǎo} — He's very skilled.
  \item \emph{Recall:} say the Mandarin for what someone is capable of, then the Mandarin for an expert, then say \emph{jìshù} again
\end{itemize}

\begin{tcolorbox}[breakable,colback=teal!4,colframe=teal!35!black,title={Before you move on}]
What does jìshù mean? (\textquotedblleft{}Skill; technique\textquotedblright{}.) Say it once more.
\end{tcolorbox}
\section[jìnbù]{jìnbù — to improve; to make progress}

\label{lesson:ZH-C260-jinbu}



\begin{quote}
Before the new one: say the Mandarin for an expert, then the Mandarin for technique.
\end{quote}

\subsection*{You'll want to know: jìnbù}
\emph{jìnbù} — \textquotedblleft{}to improve; to make progress\textquotedblright{}.

\begin{quote}
\emph{nǐ de Hànyǔ jìnbù de hěn kuài} — Your Chinese is improving fast.
\end{quote}

Its characters come later; for now it is a word for the ear.

\begin{cousinweb}
\emph{jìnbù} is built from \zh{进} 'to go forward' and \emph{bù} 'a step', both inherited from Middle Chinese, \emph{tsinH} and \emph{buH} (Guangyun, in Baxter's transcription), which glosses the first 'forward' and the second 'a pace'.
\end{cousinweb}

\begin{grammarlens}[title={What you've built}]
One more word for what I have and can do.
\end{grammarlens}

\subsection*{Your turn}
\begin{itemize}
\raggedright
  \item \emph{Say it:} \emph{jìnbù}
  \item \emph{Say it:} \emph{jìnbù}, once more
  \item \emph{Say it:} \emph{nǐ de Hànyǔ jìnbù de hěn kuài} — Your Chinese is improving fast.
  \item \emph{Recall:} say the Mandarin for an expert, then the Mandarin for technique, then say \emph{jìnbù} again
\end{itemize}

\begin{tcolorbox}[breakable,colback=teal!4,colframe=teal!35!black,title={Before you move on}]
What does jìnbù mean? (\textquotedblleft{}To improve; to make progress\textquotedblright{}.) Say it once more.
\end{tcolorbox}
